\documentclass[11pt]{article}
\usepackage{amsmath,amssymb,amsthm}
\usepackage[margin=1in]{geometry}

\newtheorem{theorem}{Theorem}
\newtheorem{corollary}{Corollary}

\title{L1 --- Equivalence Completion (Paper 7):\\
Anomaly Freedom $\Longleftrightarrow$ Yukawa Realizability under Minimal Chiral Load}
\author{Amos Jay Maley}
\date{}

\begin{document}
\maketitle

\section*{0. Scope and Non-Negotiables}
This paper is eliminative. It introduces no new primitives, no new ontology, no new anchors, and no new tensor load.
All results are invariant under admissible redescription.

We work within two admissible Standard-Model skins:
(i) a chiral gauge (unbroken-phase) bookkeeping skin and
(ii) a broken-phase (mass-eigenstate) bookkeeping skin.
The declared load is fixed: the minimal chiral fermion content and a single electroweak Higgs doublet.
No compensating sectors, additional multiplets, or anomaly-cancellation patches are permitted.

\section*{1. Apparent Freedom}
In Skin A (chiral gauge bookkeeping), anomaly constraints appear to restrict only representation/charge assignments,
while Yukawa structure is treated as modular additional data.
In Skin B (broken-phase bookkeeping), masses and mixings are encoded in Yukawa matrices and appear continuously specifiable,
with anomaly freedom treated as already handled upstream.
This motivates an \emph{apparent independence}: (charges/anomalies) $\perp$ (Yukawas/masses).

\section*{2. Admissibility Envelope}
Allowed redescriptions:
\begin{itemize}
\item gauge-basis $\leftrightarrow$ mass-basis unitary field redefinitions preserving the gauge-covariant kinetic structure,
\item unbroken $\leftrightarrow$ broken electroweak-phase bookkeeping,
\item reparameterizations preserving invariant content (phase reshufflings, equivalent parameterizations).
\end{itemize}

Standing fails (fail-closed) if any choice is non-transportable between these skins without importing structure.

\section*{3. Cross-Skin Transport: Operator Existence as an Invariant}
A Yukawa term in Skin B corresponds to a gauge-invariant operator in Skin A.
Conversely, if a gauge-invariant Yukawa operator does not exist in Skin A,
then no admissible redescription can create it in Skin B without adding new load or altering anchors.

Therefore, \emph{complete Yukawa realizability} is a transport-invariant requirement:
it is not optional data, but a constraint on the charge/representation bookkeeping.

\section*{4. Failure of Independence (Transport Contradiction)}
Assume charges/representations are chosen independently of Yukawa realizability, subject only to anomaly cancellation in Skin A.
Transport to Skin B.

If the chosen charges do not permit the full set of gauge-invariant Yukawa operators needed to pair each chiral fermion
into a mass term after symmetry breaking, then Skin B contains massless chiral remnants.
Such remnants cannot be removed by admissible field redefinition.

Eliminating them requires one of:
(i) additional Higgs multiplets,
(ii) additional fermions (including vectorlike pairs used as compensators), or
(iii) non-gauge-invariant effective operators.
Each is disallowed: (i) and (ii) add tensor load; (iii) is an external patch.
Hence independence produces standing/closure failure.

Conversely, assume Yukawa structure is freely specifiable independently in Skin B.
Transport back to Skin A.
A generic Yukawa choice presupposes the existence of corresponding gauge-invariant operators.
If charges are also freely specifiable, generically those operators do not exist.
Thus ``free Yukawas'' is not an admissible independent degree of freedom under transport.

\section*{5. Forced Tensor Relation (Core Result)}
\begin{theorem}[Anomaly--Yukawa Equivalence under Minimal Chiral Load]
Fix the minimal chiral fermion load and a single electroweak Higgs doublet in the Standard-Model skins.
Under admissible redescription and under the global constraints (no new load, no patches),
the following are equivalent tensor requirements:
\begin{enumerate}
\item \textbf{Anomaly-freedom:} all gauge and mixed anomalies vanish (the chiral-gauge skin is admissible).
\item \textbf{Complete Yukawa realizability:} there exists a complete set of gauge-invariant Yukawa operators
that generates masses for all chiral fermions in the broken-phase skin without introducing any new multiplets
or non-gauge-invariant effective terms.
\end{enumerate}
Equivalently, anomaly cancellation and full Yukawa operator realizability are the same tensor constraint expressed in different skins.
\end{theorem}

\begin{corollary}[Parameter Collapse]
Within the fixed declared load, requiring complete Yukawa realizability collapses the apparent charge freedom:
the admissible U(1) charge assignments reduce to at most an overall normalization
(i.e., a pure representational scale choice), with no additional independent parameters compatible with closure.
\end{corollary}

\begin{corollary}[Uniqueness Upgrade]
Once the equivalence is enforced, the minimal chiral gauge bookkeeping is closure-unique up to admissible redescription.
Any purported alternative that preserves the declared load but breaks the equivalence
either (i) leaves massless chiral remnants (non-transportable) or (ii) requires compensating additions or patches (forbidden).
\end{corollary}

\section*{6. Closure Statement}
The apparent modularity (anomaly constraints vs Yukawa/mass data) is a skin illusion.
Cross-skin admissibility forces their identification as a single tensor constraint.
No patching is required.

\end{document}

